% example, page 02: a tensor diagram with the styles of tensor.tex -- the expansion
% psi(r1, r2) = sum_ij C_ij phi_i(r1) phi_j(r2): the coefficient array C (square, plain
% legs i, j) on the basis functions phi (circles), which carry the positions (wavy legs).
\documentclass[border=4pt]{standalone}
\input{preamble}
\input{tensor}

\begin{document}
\begin{tikzpicture}
  \node[fnwide] (psi) at (0,0) {$\psi$};
  \draw[cont] ([xshift=-4.5mm]psi.south) -- ++(0,-0.9) node[leg, below] {$\mathbf r_1$};
  \draw[cont] ([xshift= 4.5mm]psi.south) -- ++(0,-0.9) node[leg, below] {$\mathbf r_2$};
  \node at (1.3,-0.2) {$=$};
  \node[coefwide, minimum width=22mm] (C) at (3.0,1.3) {$C$};
  \node[fn] (p1) at (2.25,0) {$\varphi$};
  \node[fn] (p2) at (3.75,0) {$\varphi$};
  \draw[disc] (C.south -| p1) -- node[leg, left] {$i$} (p1);
  \draw[disc] (C.south -| p2) -- node[leg, right] {$j$} (p2);
  \draw[cont] (p1.south) -- ++(0,-0.9) node[leg, below] {$\mathbf r_1$};
  \draw[cont] (p2.south) -- ++(0,-0.9) node[leg, below] {$\mathbf r_2$};
\end{tikzpicture}
\end{document}
